\documentclass[twocolumn]{article}
\pdfinfoomitdate=1
\pdftrailerid{}
\pdfsuppressptexinfo=-1
\usepackage{graphicx}
\usepackage{amsmath}
\usepackage{amsthm}
\usepackage{algorithm}
\usepackage{algorithmic}
\newtheorem{theorem}{Theorem}
\newtheorem{lemma}{Lemma}
\title{A tiny two-column paper, written for the PDF inventory}
\author{The Wayline Suite}
\date{}
\begin{document}
\maketitle

\begin{abstract}
This paper exists so that a PDF inventory has something to count. It has two columns,
three numbered equations, two figures, one ruled table, one algorithm, a theorem, a lemma,
an appendix, and two references.
\end{abstract}

\section{Introduction}
A tiny paper needs a tiny introduction. The method of Section~\ref{sec:method} is one
line, and Figure~\ref{fig:signal} shows it doing nothing in particular. Earlier work on
nothing includes Knuth~\cite{knuth1984} and Lamport~\cite{lamport1994}.

\subsection{Scope}
Nothing at all is in scope, which is the point. The two columns of this page are there to
see whether a converter reads the left column to the bottom before it starts on the right.

\section{Method}
\label{sec:method}
The method is one numbered equation,
\begin{equation}
  y = m x + b ,
  \label{eq:line}
\end{equation}
followed by a display that carries no number,
\[
  \int_0^1 x^2 \, dx = \frac{1}{3} ,
\]
and an aligned pair that is numbered twice:
\begin{align}
  a &= \sum_{i=1}^{n} b_i , \label{eq:sum} \\
  c &= \sqrt{a^2 + d^2} . \label{eq:root}
\end{align}

\begin{theorem}[Tininess]
\label{thm:tiny}
Every tiny paper is tiny.
\end{theorem}

\begin{proof}
By inspection of this page.
\end{proof}

\begin{lemma}
\label{lem:tinier}
Some tiny papers are tinier than others.
\end{lemma}

\begin{algorithm}
\caption{Do nothing, carefully}
\label{alg:nothing}
\begin{algorithmic}
\STATE $x \gets 0$
\FOR{$i = 1$ to $n$}
  \STATE $x \gets x + 0$
\ENDFOR
\RETURN $x$
\end{algorithmic}
\end{algorithm}

\section{Results}
Table~\ref{tab:numbers} lists the numbers the method produced, and
Figure~\ref{fig:panels} puts them side by side.

\begin{figure}
  \centering
  \includegraphics[width=0.9\columnwidth]{signal.png}
  \caption{A signal that rises from left to right and means nothing.}
  \label{fig:signal}
\end{figure}

\begin{table}
  \centering
  \caption{Numbers the method produced.}
  \label{tab:numbers}
  \begin{tabular}{|l|r|r|}
    \hline
    Run & Input & Output \\
    \hline
    first & 1 & 0 \\
    second & 2 & 0 \\
    third & 3 & 0 \\
    \hline
  \end{tabular}
\end{table}

\begin{figure*}
  \centering
  \includegraphics[width=0.8\textwidth]{panels.png}
  \caption{Two panels that span both columns, so the caption sits under the whole page width.}
  \label{fig:panels}
\end{figure*}

\appendix
\section{Proof of the lemma}
\label{app:proof}
The proof of Lemma~\ref{lem:tinier} is left to the reader, who is the point of the inventory.

\begin{thebibliography}{2}
\bibitem{knuth1984} D. E. Knuth. The \TeX book. Addison-Wesley, 1984.
\bibitem{lamport1994} L. Lamport. \LaTeX: A Document Preparation System. Addison-Wesley, 1994.
\end{thebibliography}
\end{document}
